\section{Agentic benchmarks：评价模型说什么到做什么}

Agent = language model + agent scaffold。Scaffold 包括如何调用工具、如何计划、是否维护 todo、是否写文件、是否调用子代理、如何管理长上下文。Agent benchmarks 评价的是模型和 scaffold 的组合系统。

\begin{knowledgebox}{课堂提示：agent benchmark 的归因问题}
老师在这里把评价对象从 LM 扩展成系统：同一个模型配不同 scaffold，成绩可能差很多。因此 agent benchmark 的结果不能简单归因给 base model。它测的是 model + tools + memory + planning + execution loop 的整体。
\end{knowledgebox}

\begin{figure}[H]
\centering
\includegraphics[width=0.9\textwidth]{swebench.png}
\caption{SWE-Bench：给 codebase 和 issue description，要求提交 PR，用 unit tests 评估。}
\figsource{CS336 Lecture 12 官方源引用图。}
\end{figure}

\begin{knowledgebox}{读图：SWE-Bench 为什么是 agent benchmark}
模型必须读代码、定位 bug、修改文件、运行测试并迭代。最终 metric 是 unit tests，而不是文本回答质量。它评价的是工程行动能力。图中任务结构提醒我们：评测不再是一次 forward pass，而是一个多步闭环。
\end{knowledgebox}

TerminalBench 把环境进一步抽象成通用终端。CyBench 加入安全/CTF 问题。MLE-Bench 则把数据处理、训练模型和提交结果纳入任务。这三类 benchmark 都在扩大“能力表面”：模型不只是回答，还要操作环境。

为了避免把这组图读成“又三个榜单”，可以把它们看成三种行动环境：TerminalBench 测通用终端工作流，CyBench 测安全挑战和探索，MLE-Bench 测机器学习工程项目。它们共同的问题是：模型必须在多步过程中观察、行动、失败、修正，而不是一次性输出一句答案。

\begin{figure}[H]
\centering
\includegraphics[width=0.86\textwidth]{terminal-bench.png}
\caption{TerminalBench：用通用 computer terminal environments 评价长期工具使用任务。}
\figsource{CS336 Lecture 12 官方源引用图。}
\end{figure}

\begin{figure}[H]
\centering
\includegraphics[width=0.8\textwidth]{terminal-bench-human-time.png}
\caption{TerminalBench human time：任务对人类需要的时间提供难度参考。}
\figsource{CS336 Lecture 12 官方源引用图。}
\end{figure}

\begin{figure}[H]
\centering
\includegraphics[width=0.82\textwidth]{terminal-bench-results.png}
\caption{TerminalBench results：展示 agents 在终端任务上的成功率。}
\figsource{CS336 Lecture 12 官方源引用图。}
\end{figure}

\begin{knowledgebox}{读图：TerminalBench 的通用性}
终端环境简单、通用、可脚本化，能覆盖安装、调试、运行、文件处理等真实工作流。但它也强依赖 scaffold：同一模型配不同工具循环，成绩可能显著不同。human time 图还告诉我们，难度不是抽象标签，而可以用人类完成时间校准。
\end{knowledgebox}

TerminalBench 的三张图分别承担不同角色：环境图说明任务长什么样，人类时间图给出难度标尺，结果图展示模型系统当前能做到什么。只有三者一起看，才不会把低分误解为“模型完全不会”，也不会把高分误解为“模型已经能独立工作”。评测中间还隔着 scaffold、工具权限、时间预算和错误恢复策略。

\begin{figure}[H]
\centering
\includegraphics[width=0.86\textwidth]{cybench.png}
\caption{CyBench：40 个 Capture the Flag 任务，用 first-solve time 表示难度。}
\figsource{CS336 Lecture 12 官方源引用图。}
\end{figure}

\begin{figure}[H]
\centering
\includegraphics[width=0.86\textwidth]{cybench-agent.png}
\caption{CyBench agent setup：模型需要在 CTF 环境中探索、运行工具并解题。}
\figsource{CS336 Lecture 12 官方源引用图。}
\end{figure}

\begin{figure}[H]
\centering
\includegraphics[width=0.72\textwidth]{cybench-results.png}
\caption{CyBench results：展示 agent 在安全挑战任务中的表现。}
\figsource{CS336 Lecture 12 官方源引用图。}
\end{figure}

\begin{warningbox}{读图：CyBench 的 dual-use 含义}
CTF 能测试安全推理和工具使用，也具有 dual-use 性质。模型越会做 cyber tasks，越需要同时评估防御和滥用风险。这里的“能力提升”本身不是单向好事。
\end{warningbox}

CyBench 和 MLE-Bench 进一步说明 agentic evaluation 必须报告过程，而不是只报告最终分数。CTF 任务里，模型可能靠枚举、工具调用或漏洞知识推进；Kaggle 任务里，模型可能输在环境配置、数据处理或提交策略。若没有轨迹分析，benchmark 分数很难告诉我们应该改模型、改 scaffold，还是改工具。

\begin{figure}[H]
\centering
\includegraphics[width=0.9\textwidth]{mlebench.png}
\caption{MLE-Bench：Kaggle competitions，要求处理数据、训练模型和提交结果。}
\figsource{CS336 Lecture 12 官方源引用图。}
\end{figure}

\begin{figure}[H]
\centering
\includegraphics[width=0.78\textwidth]{mlebench-results.png}
\caption{MLE-Bench results：模型在机器学习工程竞赛任务上的表现。}
\figsource{CS336 Lecture 12 官方源引用图。}
\end{figure}

\begin{knowledgebox}{读图：MLE-Bench 为什么难}
这类任务跨越数据清洗、建模、调参、提交格式、错误分析和计算预算。它更像真实 ML 工程，而不是单步问答。成绩低不一定说明模型不会 ML 知识，也可能说明 scaffold、计算预算、调试循环和提交策略不足。
\end{knowledgebox}

\begin{figure}[H]
\centering
\includegraphics[width=0.62\textwidth]{agents-overview.png}
\caption{Agent scaffold overview：planning、delegation、persistent memory、context engineering 等增强能力面。}
\figsource{Deep agents overview，本地化后供 CS336 Lecture 12 使用。}
\end{figure}

\begin{importantbox}{Agentic evaluation 的关键}
评估 agents = 评估 language model + scaffold。若不固定 scaffold，就不能把结果归因到模型本身；若固定 scaffold，又可能低估真实产品系统的能力。一个严谨 agent benchmark 应报告模型、工具、prompt、记忆、规划策略、时间预算和失败恢复规则。
\end{importantbox}

\subsection{本章小结}

Agentic benchmarks 从“说得好”转向“做得成”。它们更真实，也更难设计 valid metrics，因为 scaffold、工具、时间、环境和测试质量都会影响结果。

